% Part: computability
% Chapter: computability-theory
% Section: introduction

\documentclass[../../../include/open-logic-section]{subfiles}

\begin{document}

\olfileid[hi]{cmp}{thy}{int}
\olsection{परिचय}

तर्कशास्त्र की वह शाखा जिसे \emph{संगणनीयता सिद्धांत} कहा जाता है,
फलनों और समुच्चयों की संगणनीयता अथवा सापेक्ष संगणनीयता से जुड़े
प्रश्नों का अध्ययन करती है। क्लीन के प्रभाव का प्रमाण यह है कि इस
विषय को पहले \emph{पुनरावर्तन सिद्धांत} कहा जाता था; आज दोनों नाम
सामान्यतः प्रयुक्त होते हैं।

किसी फलन~$f\colon \Nat \pto \Nat$ को हम \emph{आंशिक संगणनीय} कहेंगे,
यदि उसे संगणना के किसी मॉडल में संगणित किया जा सकता हो। यदि $f$ पूर्ण
है, तो हम केवल यह कहेंगे कि $f$~\emph{संगणनीय} है। जिस संबंध~$R$ का
अभिलाक्षणिक फलन~$\Char{R}$ संगणनीय हो, उसे भी संगणनीय कहते हैं। यदि
$f$ और~$g$ आंशिक फलन हैं, तो $f(x) \fdefined$ लिखने का अर्थ होगा कि
$f$,~$x$ पर परिभाषित है, अर्थात् $x$,~$f$ के प्रांत में है; और
$f(x) \fundefined$ का अर्थ इसका उलटा होगा, अर्थात् $f$,~$x$ पर
परिभाषित नहीं है। $f(x) \simeq g(x)$ से हमारा आशय होगा कि या तो
$f(x)$ और $g(x)$ दोनों अपरिभाषित हैं, अथवा दोनों परिभाषित और समान हैं।

संगणना के किसी विशिष्ट मॉडल का उल्लेख किए बिना भी संगणनीयता सिद्धांत
का अध्ययन किया जा सकता है। इसके लिए यह दिखाया जाता है कि एक सार्वत्रिक
आंशिक संगणनीय फलन~$\fn{Un}(k, x)$ अस्तित्व में है। इससे आंशिक संगणनीय
फलनों का परिगणन किया जा सकता है। हम संकेत~$\cfind{k}$ अपनाएँगे, जो
$k$-वें एकस्थानी आंशिक संगणनीय फलन को निरूपित करेगा और जिसकी परिभाषा
$\cfind{k}(x) \simeq \fn{Un}(k, x)$ है। (क्लीन ने इस प्रयोजन के लिए
$\{ k \}$ का उपयोग किया था, किंतु हाल में यह संकेत उतना प्रचलित नहीं
रहा है।) कुछ अधिक सामान्य रूप में, हम मनमानी स्थान-संख्या वाले आंशिक
संगणनीय फलनों का एकसमान परिगणन कर सकते हैं; हम $\cfind{k}[n]$ लिखेंगे,
जो $k$-वें $n$-स्थानी आंशिक पुनरावर्ती फलन को निरूपित करेगा।

यदि $f(\vec x, y)$ कोई पूर्ण अथवा आंशिक फलन है, तो $\umin{y}{f
(\vec x, y)}$,~$\vec x$ का वह फलन है जो वह लघुतम~$y$ लौटाता है जिसके
लिए $f(\vec x, y) = 0$, बशर्ते $f(\vec x, 0)$, \dots,
$f(\vec x, y-1)$ सभी परिभाषित हों; यदि ऐसा कोई $y$ न हो, तो
$\umin{y}{f (\vec x, y)}$ अपरिभाषित है। यदि $R(\vec x, y)$ कोई संबंध
है, तो $\umin{y}{R(\vec x, y)}$ को उस लघुतम~$y$ के रूप में परिभाषित
किया जाता है जिसके लिए $R(\vec x, y)$ सत्य हो; दूसरे शब्दों में, वह
लघुतम~$y$ जिसके लिए $1 \tsub \Char{R}(\vec x, y) = 0$।

किसी फलन को संगणनीय दिखाने के लिए दो प्रकार से आगे बढ़ा जा सकता है:
\begin{enumerate}
\item कठोर रूप से: किसी ट्यूरिंग मशीन अथवा आंशिक पुनरावर्ती फलन का
  स्पष्ट वर्णन करें और दिखाएँ कि वह इच्छित फलन को संगणित करता है;
\item अनौपचारिक रूप से: उसे संगणित करने वाली किसी कलनविधि का वर्णन
  करें और चर्च की प्रस्थापना का सहारा लें।
\end{enumerate}
इन दोनों के बीच कोई तीखी सीमा नहीं है; किसी कलनविधि के विस्तृत वर्णन
में इतनी सूचना होनी चाहिए कि यह अपेक्षाकृत स्पष्ट हो जाए कि सिद्धांततः
उचित ट्यूरिंग मशीन अथवा आंशिक पुनरावर्ती परिभाषाओं का अनुक्रम कैसे
बनाया जा सकता है। पूर्णतः कठोर परिभाषाएँ प्रायः विशेष ज्ञानवर्धक नहीं
होंगी, इसलिए हम इन दोनों दृष्टियों के बीच उपयुक्त संतुलन खोजेंगे; संक्षेप
में, ऐसे सहज किंतु कठोर प्रमाण खोजेंगे जिनसे सटीक परिभाषाएँ प्राप्त की
जा सकें।

\end{document}
